\chapter{Mở Lại Không Khí}
\chapterintro{Ba phút đầu đủ cứu cả tiết}{Khi lớp đang nguội, điều quan trọng không phải là nói nhiều hơn mà là đổi nhịp đúng chỗ.}

\section{Ba từ kéo cả lớp ngẩng lên}
\begin{tinhhuong}{Tình huống}
Lớp vừa vào tiết nhưng mắt học sinh còn đang ở ngoài hành lang. Thầy cô nói tiếp thì lớp nghe rất ít, nhắc trật tự thì không khí lại nặng.
\end{tinhhuong}
\begin{goiydungngay}
Mở bằng ba từ ngắn và rõ: \textit{``Ngẩng lên nào.''} Sau đó im hai giây. Chỉ nói tiếp khi phần lớn ánh mắt đã quay lại.
\begin{itemize}
\item Không kéo dài thành bài nhắc nhở.
\item Giọng nhẹ nhưng dứt.
\item Nếu lớp còn lơ, lặp lại đúng một lần.
\end{itemize}
\end{goiydungngay}
\begin{cauchot}
Mở lớp bằng sự rõ ràng, không mở lớp bằng càm ràm.
\end{cauchot}
\ngancach

\section{Đố nhanh một câu không cần đúng}
\begin{tinhhuong}{Tình huống}
Đầu giờ lớp mệt, hỏi kiến thức cũ thì ít em dám trả lời vì sợ sai. Cả lớp chờ người khác lên tiếng trước.
\end{tinhhuong}
\begin{goiydungngay}
Đưa một câu đố nhanh kiểu: \textit{``Nếu bài học này là một món ăn, nó sẽ là món gì?''} hoặc \textit{``Nếu chủ đề hôm nay có màu, em chọn màu nào?''}
\begin{itemize}
\item Nói rõ: không có đáp án đúng sai.
\item Gọi 3 em trả lời thật nhanh.
\item Từ câu trả lời vui, nối ngay sang bài chính.
\end{itemize}
\end{goiydungngay}
\begin{cauchot}
Khi học sinh hết sợ sai, lớp bắt đầu sống lại.
\end{cauchot}
\ngancach

\section{Đổi vai trong 30 giây}
\begin{tinhhuong}{Tình huống}
Lớp chậm vào bài vì các em đang ở trạng thái nghe thụ động. Cần một động tác ngắn để học sinh chuyển sang chủ động.
\end{tinhhuong}
\begin{goiydungngay}
Nói: \textit{``30 giây tới, em là thầy cô. Em sẽ nói gì để lớp chú ý?''}
\begin{itemize}
\item Cho 2 em nói liền nhau.
\item Nhặt một câu hay nhất rồi dùng luôn làm câu mở tiết.
\item Cả lớp thường bật cười, nhưng cũng nhập cuộc rất nhanh.
\end{itemize}
\end{goiydungngay}
\begin{cauchot}
Cho học sinh đổi vai ngắn là cách kéo các em trở lại vai người học thật.
\end{cauchot}
\ngancach

\section{Một vật trên bàn cũng thành câu mở tiết}
\begin{tinhhuong}{Tình huống}
Thầy cô muốn gọi lớp vào bài nhanh nhưng không muốn mở bằng một câu hỏi quá trừu tượng.
\end{tinhhuong}
\begin{goiydungngay}
Chỉ vào một vật gần nhất như cây bút, quyển sách, chai nước rồi hỏi: \textit{``Nếu vật này phải đại diện cho tiết học hôm nay, nó đang nhắc mình điều gì?''}
\begin{itemize}
\item Chọn vật thật gần để học sinh không cần nghĩ xa.
\item Mời 2 đến 3 câu trả lời thật ngắn.
\item Từ đó nối luôn sang nội dung bài.
\end{itemize}
\end{goiydungngay}
\begin{cauchot}
Điều gần nhất trước mắt thường là chiếc cầu nhanh nhất để kéo lớp vào bài.
\end{cauchot}
\ngancach

\section{Cho lớp tự đo nhiệt độ đầu giờ}
\begin{tinhhuong}{Tình huống}
Đầu tiết học sinh chưa hẳn ồn, nhưng khí lớp đang nguội và rời rạc.
\end{tinhhuong}
\begin{goiydungngay}
Hỏi: \textit{``Nếu chấm trạng thái lớp mình từ 1 đến 5, giờ đang ở mức mấy?''}
\begin{itemize}
\item Cho giơ tay hoặc giơ ngón tay thật nhanh.
\item Không phân tích dài, chỉ đọc bầu không khí lớp.
\item Chốt bằng câu: \textit{``Rồi, mình cùng kéo từ mức này lên cao hơn một nấc.''}
\end{itemize}
\end{goiydungngay}
\begin{cauchot}
Khi lớp tự nhận ra mình đang ở đâu, các em dễ cùng thầy cô đi tiếp hơn.
\end{cauchot}
